\documentclass[12pt]{article}

% ── Packages ────────────────────────────────────────────────────────────────
\usepackage[margin=1.25in]{geometry}
\usepackage{setspace}
\usepackage[utf8]{inputenc}
\usepackage[T1]{fontenc}
\usepackage{microtype}
\usepackage{amsmath, amssymb}
\usepackage{booktabs}
\usepackage{graphicx}
\usepackage{caption}
\usepackage{subcaption}
\usepackage{float}
\usepackage{hyperref}
\usepackage{natbib}
\usepackage{ragged2e}
\usepackage{pgfmath}   % for \round and \formatInt helpers below
 \usepackage{clipboard}
 \newclipboard{clipboardForResponses}

% ── Number-formatting helpers (same as in the real paper) ───────────────────
\newcommand{\round}[1]{%
  \pgfmathparse{#1}%
  \pgfmathprintnumber[fixed, fixed zerofill, precision=2]{\pgfmathresult}%
}
\newcommand{\formatInt}[1]{%
  \pgfkeys{/pgf/number format/.cd,
    fixed, precision=0, 1000 sep={,}}%
  \pgfmathprintnumber{#1}%
}

% ── Load all scalars produced by the analysis code ──────────────────────────
%    refresh_output.sh copies output/scalars.tex → writing/paper/output/scalars.tex
\newcommand{\wonderfulConstant}{78.72258899555547} % written by Titl0001 running STATA on 15 Mar 2026 at 21:59:04


% ── Document metadata ───────────────────────────────────────────────────────
\onehalfspacing
\bibliographystyle{apalike}

\title{An Example Paper}
\author{John Doe}
\date{\today}

% ════════════════════════════════════════════════════════════════════════════
\begin{document}
\maketitle

% ── Abstract ────────────────────────────────────────────────────────────────
\begin{abstract}
  This example paper demonstrates how a replication package connects
  analysis code to a \LaTeX{} writeup.
  Every number in the text, every table cell, and every figure is
  generated by \texttt{code/make.sh} and copied to
  \texttt{writing/paper/output/} by \texttt{refresh\_output.sh}.
  Re-running the code and recompiling this file always produces
  an up-to-date paper with no manual editing. This is a wonderful constant: $\wonderfulConstant$ directly from Stata.
\end{abstract}

\newpage

% ════════════════════════════════════════════════════════════════════════════
\section{Introduction}
\label{sec:intro}

The introduction of wonderful paper. And I again repeat the wonderful constant: $\wonderfulConstant$.


% ════════════════════════════════════════════════════════════════════════════
\section{Data}
\label{sec:data}

Table~\ref{tab:summary} reports summary statistics.
% ── TABLE: summary statistics ───────────────────────────────────────────────
\begin{table}[H]
  \centering
  \caption{Summary Statistics}
  \label{tab:summary}
  {
\def\sym#1{\ifmmode^{#1}\else\(^{#1}\)\fi}
\begin{tabular}{l*{1}{ccccc}}
\hline\hline
                    &\multicolumn{5}{c}{}                                            \\
                    &       count&        mean&          sd&         min&         max\\
\hline
tender\_estimated\_price&     1567751&    1.25e+08&    2.05e+09&           1&    2.25e+12\\
\hline\hline
\end{tabular}
}

  \begin{flushleft}
    \footnotesize
    \textit{Notes:}
    The table reports means, medians and standard deviations for the
    key variables.
  \end{flushleft}
\end{table}


% ════════════════════════════════════════════════════════════════════════════
\section{Results}
\label{sec:results}
\subsection{Parallel Trends Validation}
\label{sec:pretrends}
Before turning to the treatment effect estimates, we examine whether 
treated and not-yet-treated municipalities followed similar emission 
trajectories prior to threshold crossing. Panel A of Table 
\ref{tab:dynamic} and Figure \ref{fig:co2pl_lev5_10} report pre-treatment 
coefficients from the event study specification over the window 
$T_{-5}$ to $T_{+5}$.

The pre-treatment coefficients are small and close to zero throughout, 
with point estimates ranging from $-0.000109$ to $-0.000268$. All 
four periods from $T_{-5}$ through $T_{-2}$ are clearly insignificant 
($p > 0.37$). The coefficient at $T_{-1}$ is $-0.000268$ 
(SE $= 0.000153$, $p = 0.080$). While marginally 
significant in isolation, the joint test of all five pre-treatment coefficients  yields $\chi^2(5) = 3.80$ ($p = 0.579$), and the 
average pre-treatment coefficient is $-0.000177$ 
(SE $= 0.000109$, $p = 0.104$). Both confirm that 
this pattern reflects sampling variation rather than a systematic 
pre-existing differential trend. Moreover, the 
magnitude is small --- less than one sixth of the 
average post-treatment effect --- providing no 
economically meaningful evidence of pre-treatment 
divergence. We therefore proceed on the basis that 
the conditional parallel trends assumption holds.
\label{highlightedResultInPaperPage}
\Copy{highlightedResultInPaper}{Here is a new paragraph that demonstrates how not copy paste to the answer to reviewers, but I have it automatted}.

% ════════════════════════════════════════════════════════════════════════════
\section{Conclusion}
\label{sec:conclusion}

This paper demonstrates the ideal workflow for a replication package.
% ════════════════════════════════════════════════════════════════════════════
\bibliography{references}

\end{document}
